\begin{abstract}

We present \method,
a feed-forward model that directly predicts all articulations of an object,
including part segmentation, kinematic structure, motion type, and continuous joint parameters,
from two distinct states of the same object.
Unlike existing \emph{single}-state feed-forward methods,
\method exploits the observed state change as direct evidence of articulation,
by learning to match the corresponding motion across the \emph{two} states,
rather than relying heavily on category-level priors learned from limited articulated-object annotations.
\method is a Siamese motion-based articulation transformer,
along with a cycle-consistency training objective that encourages consistent predictions under reversed state order.
We demonstrate the importance of the two-state formulation and cycle-consistency training for learning feed-forward articulations through extensive experiments.
We also show that \method achieves state-of-the-art performance on articulated structure recovery, with better generalization to novel object categories and more complex kinematic layouts.
% Additionally, we show that, with the help of VLMs, even a single state input can be extended to multiple states as conditioning for our model, further improving its performance and generalization.

% Recovering articulated 3D objects requires understanding not only their geometry, but also how their parts move and constrain one another.
% Existing feed-forward methods usually infer this structure from a single static observation, forcing the model to rely heavily on category-level priors learned from limited articulated-object annotations.
% We instead study articulation inference from two distinct states of the same object, where the observed state change directly exposes the moving parts and their underlying kinematics.
% We introduce \method, a feed-forward framework that takes two unorganized point clouds as input and predicts part segmentation, kinematic structure, motion type, and continuous joint parameters.
% \method uses a Siamese motion-based articulation transformer with shared part queries, query-to-point masked attention, point-to-query attention, and a cycle-consistency objective that encourages consistent predictions under reversed state order.
% This design allows the model to exploit cross-state geometric evidence without requiring point correspondences between the two observations, while retaining fast single-pass inference at test time.
% Experiments on articulated object benchmarks show that the two-state formulation substantially improves over single-state prediction and achieves state-of-the-art performance on articulated structure recovery, with ablations validating the contribution of cross-state reasoning and consistency training.
\end{abstract}
